% 使用 XeLaTeX 编译；中文字体使用 TeX Live 自带的 Fandol 字体。
\documentclass[11pt,a4paper,fontset=fandol]{ctexart}
\usepackage[margin=2.1cm]{geometry}
\usepackage{amsmath,amssymb,mathtools}
\usepackage{array,booktabs}
\usepackage{fancyhdr}
\usepackage[hidelinks]{hyperref}

\setlength{\parindent}{0pt}
\setlength{\parskip}{0.3em}
\setlength{\headheight}{15pt}
\renewcommand{\arraystretch}{1.05}
\pagestyle{fancy}
\fancyhf{}
\fancyhead[L]{MATH1205\quad 线性代数}
\fancyhead[R]{第四次作业}
\fancyfoot[C]{\thepage}
\renewcommand{\headrulewidth}{0.3pt}
\ctexset{section={format=\large\bfseries,beforeskip=1.1ex,afterskip=0.8ex},
  subsection={format=\normalsize\bfseries,beforeskip=0.9ex,afterskip=0.5ex}}
\newcommand{\T}{^{\mathsf T}}
\newcommand{\ans}[1]{\boxed{#1}}

\begin{document}
\begin{center}
  {\LARGE\bfseries 线性代数第四次作业}\par
  \vspace{0.25em}
  {\small MATH1205}
\end{center}
\vspace{-0.5em}

\section*{第 1 题（原题号 20）：用行列式的性质计算}
以下用 $R_i,C_j$ 分别表示第 $i$ 行、第 $j$ 列。

\subsection*{(1)}
\[
D=\begin{vmatrix}
x^2+1&yx&zx\\
xy&y^2+1&zy\\
xz&yz&z^2+1
\end{vmatrix}.
\]
记 $u=(x,y,z)\T$，$e_1,e_2,e_3$ 为三维标准单位列向量，则
\[
D=\det(e_1+xu,\ e_2+yu,\ e_3+zu).
\]
利用行列式对每一列的线性性质展开。凡含有两列或三列 $u$ 的项，
均因有相同列而为零，故只剩下
\begin{align*}
D&=\det(e_1,e_2,e_3)+x\det(u,e_2,e_3)
  +y\det(e_1,u,e_3)+z\det(e_1,e_2,u)\\
 &=1+x\cdot x+y\cdot y+z\cdot z.
\end{align*}
所以
\[
\ans{D=1+x^2+y^2+z^2}.
\]

\subsection*{(2)}
\[
D=\begin{vmatrix}
x&y&x+y\\
y&x+y&x\\
x+y&x&y
\end{vmatrix}.
\]
作列变换 $C_1\leftarrow C_1+C_2+C_3$，再从第一列提出公因子 $2(x+y)$，得
\[
D=2(x+y)\begin{vmatrix}
1&y&x+y\\
1&x+y&x\\
1&x&y
\end{vmatrix}.
\]
再作 $R_2\leftarrow R_2-R_1$、$R_3\leftarrow R_3-R_1$，按第一列展开：
\begin{align*}
D&=2(x+y)\begin{vmatrix}
1&y&x+y\\
0&x&-y\\
0&x-y&-x
\end{vmatrix}
=2(x+y)\begin{vmatrix}x&-y\\x-y&-x\end{vmatrix}\\
&=2(x+y)(-x^2+xy-y^2).
\end{align*}
因此
\[
\ans{D=-2(x+y)(x^2-xy+y^2)=-2(x^3+y^3)}.
\]
上述提取公因子的等式在 $x+y=0$ 时仍成立，无须除以 $x+y$。

\newpage
\section*{第 1 题（续）}

\subsection*{(3)}
\[
D=\begin{vmatrix}
a^2&(a+1)^2&(a+2)^2&(a+3)^2\\
b^2&(b+1)^2&(b+2)^2&(b+3)^2\\
c^2&(c+1)^2&(c+2)^2&(c+3)^2\\
d^2&(d+1)^2&(d+2)^2&(d+3)^2
\end{vmatrix}.
\]
对任意数 $t$，有
\[
(t+3)^2-3(t+2)^2+3(t+1)^2-t^2=0.
\]
因此作列变换 $C_4\leftarrow C_4-3C_3+3C_2-C_1$ 后，第四列全为零；
此变换不改变行列式的值，所以
\[
\ans{D=0}.
\]

\subsection*{(4)}
\[
D=\begin{vmatrix}
0&x&y&z\\
x&0&z&y\\
y&z&0&x\\
z&y&x&0
\end{vmatrix}.
\]
先作 $R_1\leftarrow R_1+R_2$、$R_3\leftarrow R_3+R_4$，
再作 $C_2\leftarrow C_2-C_1$、$C_4\leftarrow C_4-C_3$，得到
\[
D=\begin{vmatrix}
x&0&y+z&0\\
x&-x&z&y-z\\
y+z&0&x&0\\
z&y-z&x&-x
\end{vmatrix}.
\]
交换第 $2,3$ 行，再交换第 $2,3$ 列。两次交换的符号相消，故
\[
D=\begin{vmatrix}
x&y+z&0&0\\
y+z&x&0&0\\
x&z&-x&y-z\\
z&x&y-z&-x
\end{vmatrix}.
\]
右上角的 $2\times2$ 块全为零，按最后两列作拉普拉斯展开，得
\begin{align*}
D&=\begin{vmatrix}x&y+z\\y+z&x\end{vmatrix}
   \begin{vmatrix}-x&y-z\\y-z&-x\end{vmatrix}\\
 &=\bigl[x^2-(y+z)^2\bigr]\bigl[x^2-(y-z)^2\bigr].
\end{align*}
所以
\[
\ans{D=(x+y+z)(x-y-z)(x-y+z)(x+y-z)}.
\]

\newpage
\section*{第 2 题（原题号 21）：求解方程 $f(x)=0$}
\[
f(x)=\begin{vmatrix}
1&x&x^2&x^3\\
1&1&1&1\\
1&-1&1&-1\\
1&2&4&8
\end{vmatrix}.
\]
令 $p(t)=(t-1)(t+1)(t-2)=t^3-2t^2-t+2$。
作列变换 $C_4\leftarrow C_4-2C_3-C_2+2C_1$，最后一列成为 $(p(x),0,0,0)\T$。
沿第四列展开，并对所得三阶行列式作 $R_2\leftarrow R_2-R_1$、$R_3\leftarrow R_3-R_1$，有
\[
f(x)=-p(x)\begin{vmatrix}1&1&1\\1&-1&1\\1&2&4\end{vmatrix}
=-p(x)\begin{vmatrix}1&1&1\\0&-2&0\\0&1&3\end{vmatrix}
=6(x-1)(x+1)(x-2).
\]
所以方程的全部解为
\[
\ans{x=-1,\ 1,\ 2}.
\]

\section*{第 3 题（原题号 23）：余子式与代数余子式}
已知
\[
D=\begin{vmatrix}
3&1&0&4\\0&2&-1&1\\1&1&2&1\\3&5&2&7
\end{vmatrix},\qquad A_{ij}=(-1)^{i+j}M_{ij}.
\]
第四行各元素的余子式不依赖于第四行的元素。因此，把第四行换成
$(b_1,b_2,b_3,b_4)$ 后，新行列式沿第四行展开为 $\sum_{j=1}^4 b_jA_{4j}$。

\subsection*{(1) 求 $4M_{42}+2M_{43}+2M_{44}$}
因 $M_{42}=A_{42}$、$M_{43}=-A_{43}$、$M_{44}=A_{44}$，有
\[
4M_{42}+2M_{43}+2M_{44}
=\begin{vmatrix}
3&1&0&4\\0&2&-1&1\\1&1&2&1\\0&4&-2&2
\end{vmatrix}=\ans{0},
\]
其中最后一步使用了第四行等于第二行的 $2$ 倍。

\subsection*{(2) 求 $A_{41}+A_{42}+A_{43}+A_{44}$}
将第四行换成 $(1,1,1,1)$，再作 $R_3\leftarrow R_3-R_4$，沿第三行展开：
\begin{align*}
A_{41}+A_{42}+A_{43}+A_{44}
&=\begin{vmatrix}3&1&0&4\\0&2&-1&1\\0&0&1&0\\1&1&1&1\end{vmatrix}
=\begin{vmatrix}3&1&4\\0&2&1\\1&1&1\end{vmatrix}\\
&=3(2-1)-(0-1)+4(0-2)=\ans{-4}.
\end{align*}

\newpage
\section*{第 4 题（原题号 26）：用克拉默法则解方程组}

\subsection*{(1)}
\[
\begin{cases}
x_1+2x_2+4x_3=31,\\
5x_1+x_2+2x_3=29,\\
3x_1-x_2+x_3=10.
\end{cases}
\]
系数行列式为
\[
D=\begin{vmatrix}1&2&4\\5&1&2\\3&-1&1\end{vmatrix}
=1(1+2)-2(5-6)+4(-5-3)=-27\ne0.
\]
分别用常数列 $(31,29,10)\T$ 替换系数矩阵的第 $1,2,3$ 列：
\begin{align*}
D_1&=\begin{vmatrix}31&2&4\\29&1&2\\10&-1&1\end{vmatrix}
 =93-18-156=-81,\\[0.3em]
D_2&=\begin{vmatrix}1&31&4\\5&29&2\\3&10&1\end{vmatrix}
 =9+31-148=-108,\\[0.3em]
D_3&=\begin{vmatrix}1&2&31\\5&1&29\\3&-1&10\end{vmatrix}
 =39+74-248=-135.
\end{align*}
由克拉默法则，方程组有唯一解
\[
\ans{x_1=\frac{D_1}{D}=3,\qquad x_2=\frac{D_2}{D}=4,\qquad
x_3=\frac{D_3}{D}=5}.
\]
代入原方程，左端分别为 $3+8+20=31$、$15+4+10=29$、$9-4+5=10$，验算成立。

\newpage
\section*{第 4 题（续）}
\subsection*{(3)}
原方程组写成 $BX=b$，其中 $X=(x_1,x_2,x_3,x_4,x_5)\T$，且
\[
B=\begin{pmatrix}
3&2&0&0&0\\
1&3&2&0&0\\
0&1&3&2&0\\
0&0&1&3&2\\
0&0&0&1&3
\end{pmatrix},\qquad b=(1,0,0,0,0)\T.
\]
\textbf{第一步：计算系数行列式。}
记 $d_k$ 为主对角元全为 $3$、上副对角元全为 $2$、下副对角元全为 $1$ 的
$k$ 阶三对角行列式，并约定 $d_0=1$，则 $d_1=3$。
按第一行展开，第二项对应的余子式再沿第一列展开，得到
\[
d_k=3d_{k-1}-2d_{k-2}\qquad(k\ge2).
\]
于是
\[
\begin{array}{c|rrrrrr}
k&0&1&2&3&4&5\\\hline
d_k&1&3&7&15&31&63
\end{array}
\qquad\Longrightarrow\qquad D=d_5=63\ne0.
\]

\textbf{第二步：计算替换列后的行列式。}
用 $b$ 替换 $B$ 的第 $i$ 列得到 $D_i$。沿第 $i$ 列展开，该列只有第 $1$ 行的元素 $1$ 非零，故
\[
D_i=(-1)^{1+i}\det B_{\widehat{1},\widehat{i}},
\]
其中 $B_{\widehat{1},\widehat{i}}$ 表示从 $B$ 中删去第 $1$ 行、第 $i$ 列的矩阵。
此矩阵左上角的 $(i-1)$ 阶块为对角元全是 $1$ 的上三角矩阵，
左下角块为零，右下角是对应于 $d_{5-i}$ 的三对角矩阵。因此
\[
\det B_{\widehat{1},\widehat{i}}=d_{5-i},\qquad
D_i=(-1)^{i+1}d_{5-i}\quad (i=1,\ldots,5).
\]
这里 $i=1$ 或 $i=5$ 时出现的空块，其行列式按 $1$ 计。例如
\[
D_3=(-1)^{1+3}
\begin{vmatrix}
1&3&0&0\\
0&1&2&0\\
0&0&3&2\\
0&0&1&3
\end{vmatrix}
=\begin{vmatrix}3&2\\1&3\end{vmatrix}=7.
\]
故
\[
(D_1,D_2,D_3,D_4,D_5)=(31,-15,7,-3,1).
\]

\textbf{第三步：应用克拉默法则。}
因为 $D=63\ne0$，所以方程组有唯一解
\[
\ans{(x_1,x_2,x_3,x_4,x_5)
=\frac1{63}(31,-15,7,-3,1)}.
\]
代入验算得 $B(31,-15,7,-3,1)\T=(63,0,0,0,0)\T$，除以 $63$ 后即为 $b$。

\newpage
\section*{第 5 题（原题号 27 (1)）：齐次方程组有非零解的条件}
已知齐次方程组
\[
\begin{cases}
x_1-x_2=\lambda x_1,\\
-x_1+2x_2-x_3=\lambda x_2,\\
-x_2+x_3=\lambda x_3
\end{cases}
\]
有非零解，求 $\lambda$。

移项得
\[
\begin{pmatrix}
1-\lambda&-1&0\\
-1&2-\lambda&-1\\
0&-1&1-\lambda
\end{pmatrix}
\begin{pmatrix}x_1\\x_2\\x_3\end{pmatrix}
=\begin{pmatrix}0\\0\\0\end{pmatrix}.
\]
齐次线性方程组有非零解的充分必要条件是系数行列式为零。因此
\begin{align*}
0&=\begin{vmatrix}
1-\lambda&-1&0\\
-1&2-\lambda&-1\\
0&-1&1-\lambda
\end{vmatrix}\\
&=(1-\lambda)\bigl[(2-\lambda)(1-\lambda)-1\bigr]-(1-\lambda)\\
&=(1-\lambda)\bigl[(2-\lambda)(1-\lambda)-2\bigr]\\
&=(1-\lambda)(\lambda^2-3\lambda)
=\lambda(1-\lambda)(\lambda-3).
\end{align*}
解得
\[
\ans{\lambda=0,\ 1,\ 3}.
\]
这三个值确实都能产生非零解。例如
\[
\begin{array}{c|c}
\lambda&\text{一组对应的非零解 }(x_1,x_2,x_3)\\\hline
0&(1,1,1)\\
1&(1,0,-1)\\
3&(1,-2,1)
\end{array}
\]
分别代入原方程组即可验证，故上述参数值全部符合题意。

\end{document}
